\section{Action and Contemplation}

In the 1950s, several of \textbf{Julius Evola}'s essays were published in the English language journal, East and West, by the \textbf{Istituto Italiano per l'Africa e l'Oriente}. In January 1954, Evola's review of \textbf{Rene Guenon}'s \textit{Crisis of the Modern World} appeared. The particular topic that interests now is the notion of the elite. Guenon refers to them as the intellectual elite, although, of course, he does not mean ``intellectual'' as referring to any sort of academic discipline, and still less to the intelligentsia of writers, journalists, artists, professors, and so on. For Guenon, the intellect refers to a spiritual gnosis, a form of knowing superior to the rational or scientific mind. Evola describes it:

\begin{quotex}
Guenon does not use the expression ``intellectuality'' in its generally accepted meaning: those to whom he refers are not ``intellectuals'', but men of superior character whose formation has been on traditional lines and who possess a knowledge of metaphysics.
\end{quotex}


We won't be surprised that Evola brings up the topic of action at this point. Rather than an intellectual order, per se, Evola would prefer to see an Order, along the lines of the Templars\footnote{\url{https://www.newadvent.org/cathen/14493a.htm}}, Ismaelites\footnote{\url{https://www.onislam.net/english/ask-the-scholar/ideologies-movements-and-religions/174770.html?Religions=}}, or the Teutonic Knights\footnote{\url{https://www.newadvent.org/cathen/14541b.htm}}. He explains:

\begin{quotex}
An Order represents a superior form of life within the framework of a life of action, which may have a metaphysical and traditional ``dimension'' while at the same time remaining in a more direct touch with the world of reality and with historic facts.
\end{quotex}


So far, so good. Yet, once again, Evola muddles things a little by misunderstanding the proper relationship between contemplation and action. This bane goes on to affect his understanding of castes, the roles of priests, kings, etc. In order to be true to tradition, Evola often resorts to equivocal formulations. So while we are sympathetic to the idea of an Order acting in history, it is time to clear things up. According to Evola:

\begin{quotex} 
Guenon believes that one of the causes of the crisis of the modern world is to be found in the theoretical and practical denial of the priority that should be given to knowledge, contemplation, and pure intellectuality over action.
\end{quotex}


First of all, we can agree with Evola that by ``action'', he does not mean:

\begin{quotex}
Disorderly, unenlightened, and purposeless activity, dominated exclusively by contingent and material considerations, aiming only at worldly achievements, which is now the only form of action modern civilization recognizes and admires.
\end{quotex}


However, Evola is still not ready to conceive the priority of the intellect over action. Contemplation is the symbol of the priestly caste while action is that of the warrior caste and the king. It is true that Guenon claimed that there is a higher principle that unites the two, and this is manifested by the Priest-King of the Ancient City who embodied both principles. Evola concludes from this that neither of the two principles can claim priority over the other. This is clearly an erroneous conclusion, since the two principles need be related in some way, and is especially puzzling coming from someone who rejects any form of egalitarianism. To understand why, it is necessary to identify the underlying unifying principle, something neither Guenon nor Evola bothered to do. Let's start with schema \ref{tab:PrinciplesFaculties}.

\begin{table}[h]
\begin{tabular}{ccc}\toprule
\textbf{Principle} & \textbf{Mental Faculty} & \textbf{Transcendental}\\\midrule

Contemplation & Thinking & Truth \\

Action & Willing & Goodness \\\bottomrule
\end{tabular}
\caption{Principles and Faculties}
\label{tab:PrinciplesFaculties}
\end{table}

From this chart, we can immediately infer that the unifying principle is the One, in Neoplatonic terms, or God (Being), in medieval terms. By the metaphysical doctrine of divine simplicity, God, or the One, cannot be divided; it is nondual, hence Truth and Goodness are the same. However, as they manifest, they separate; it is there that we see their relationship. Let's go back to Evola's rejection of action as understood by the world.

\begin{description}

\item[Disorderly Activity]

The opposite of this is orderly activity, i.e., activity in harmony with the Logos, or the Cosmic Law. It is the intellect that knows the Logos.
\item[Unenlightened]

The enlightened mind is the mind that ``knows''
\item[Purposeless]

Action cannot determine its own purpose; it is the intellect that chooses it.
\item[Worldly achievements]

Of course, true action aims at supernatural achievements.
\end{description}


Clearly, the intellectual principle, in relationship to action, is the ordering principle, guiding enlightened action, providing it with a purpose and a superior aim. Another way to see this, is to use the doctrine of the symbolism of the cross, as described by Guenon. Action, in ``the world of reality'' ``historic facts'' represents the horizontal direction and the intellect represents the vertical direction.

Whatever Evola may have meant, Tradition does not accept the independence of the will in respect to the intellect. This is the false doctrine: \emph{Sic volo, sic jubeo, sit pro ratione voluntas} ``thus I wish, thus I command, my will stands in place of reason'' (Juvenal) . Since Evola has defended Reason, as long as its first principles are metaphysical doctrines, he cannot possibly mean that.

If he means that the active life in the world, if engaged in consciously, can lead to higher states of awareness, then we are in agreement. For there is a second way to understand action. Action, in the sense of bringing the potential into act, can be understood vertically. As Guenon states, it is a goal for some to actualize all their possibilities, i.e., to make real what exists only in potential. In that struggle against instinctive and inferior forces, both within and without, \emph{I simultaneously reveal myself, create myself, and know myself}.

H/T Avery


\flrightit{Posted on 2012-05-22 by Cologero}

\begin{center}* * *\end{center}

\begin{footnotesize}
\begin{sffamily}
\texttt{Charlotte on 2012-05-23 at 10:25 said:}

Ah now this is where Tomberg and Guenon do not agree, as the former gives explicit precedence to the artist over `initiated mages' of the mind. The reason being that artists have the power to inspire love, promote peace and directly convey the word of God to relatively large numbers of people -- to enliven the state and provide a channel by which divine guidance might move.

\hfill

\texttt{logres on 2012-05-23 at 18:43 said:}

Charlotte, when you say ``Art'' what are you meaning?

\hfill

\texttt{logres on 2012-05-23 at 18:49 said:}

Sic volo, sic jubeo. Sit pro ratione voluntas. Lutherus ita vult, et ait se doctorem esse super omnes doctores in toto papatu.

Luther, quoting Juvenal

Not surprising.

\hfill

\texttt{Charlotte on 2012-05-23 at 19:13 said:}

well for example Picasso's Guernica, T S Eliot's Four Quartet's, Tolkien's Lord of the Rings, Bob Marley singing One Love, Sylvie Guillem dancing... .

\hfill

\texttt{Charlotte on 2012-05-23 at 19:16 said:}

Here's another one:

There are more things in heaven and earth, Horatio, Than are dreamt of in your philosophy.

\hfill

\texttt{logres on 2012-05-23 at 22:37 said:}

The Nous manifests in Art, but also in philosophy. This doesn't even address what is un-manifest. I think what you are ``seeing'' is Tomberg's debt to the French sensibilities of the hermetic current. But I'm still not sure what you mean by Art. I can't picture Tomberg correcting Guenon's metaphysics, although it's fair to say he might find the company of Schuon more genial to his own sensibilities.

\hfill

\texttt{Charlotte on 2012-05-24 at 05:29 said:}

I never said it didn't appear in philosopy but I can quite assure you Tomberg DID say this -- Guenon isn't the last word on truth you know, he's just one more person striving to bring it to light -- and in many cases he succeeds. Tomberg also, by the way, warned vehemently against founding `schools' dedicated to following one mage or another, as they only end up in closed circles of idol worship -- something you rightly pointed out are ultimately `satanic' as they contain no possibility of freedom. In fact he pleaded with his readers not to do so.

In Tomberg's book the artist is `higher' -- more `elite', than the thinker or mere initiate, for the reasons I pointed out, which seems logical enough to me, I know very well when I'm in the presence of genius and take great joy from its manifestation... .he was humble and discerning enough to comprehend the truth and happy to consider that there were others who'd reached greater heights than he himself had achieved. It must have been a relief for someone considerd by many to be an Avatar -- a title he shunned. He was wise, he didn't want mummifying by apparently well meaning supporters -- the best school masters prefer free thinking to blind acceptance of every word they utter

Art can be appreciated by a small child and they'll be first through the doors when heaven opens -- the heart and soul can go further than the mind. Dante knew this and I presume he is considered `elite' and therefore `in the knoww', `gnostic', as it were?

I have been in the heaven which takes most of his light, and I have seen things which cannot be told, possibly, by anyone who comes down from up there.

Because, approaching the object of its desires, our intellect is so deeply absorbed that memory cannot follow it all the way.

Nevertheless, what I was able to store up of that holy kingdom, in my mind, will now be the matter of my poem.

\hfill

\texttt{logres on 2012-05-25 at 00:12 said:}

Tomberg frankly admits he is pursuing an elucidation of the personal path (hermeticism); he isn't trying to provide a metaphysic for the dying West. I don't understand why the ``flight to Art''- doesn't it have elitists and so forth as well? How is Guenon not an ``artist'' of the Intellect? Tomberg and Evola both defend the ``initation from above'', Evola in the res publica, Tomberg in private contemplation. But I've only read MotT and a few articles -- VT privileges (for purposes of his work) the French hermetic tradition, into which he ``incarnates'' for the duration of the work. I agree he is a consummate artist, as well as a better Catholic than Guenon. Are you sure that Tomberg privileges Art over Intellect? Doesn't he see them as two facets to one Truth? If he read LOTR, could we not also imagine him being interested in the thought and life of the man who created it, his letters, the Silmarillion, the diaries, the whole theology of the Inklings? I think he would.

\hfill

\texttt{logres on 2012-05-25 at 00:16 said:}

Consider above, \& see what you think... .

\hfill

\texttt{Charlotte on 2012-05-25 at 10:17 said:}

Yes I am quite sure, and it is to do with the primarcy of the heart over mind as our guide. It should be no surprise if Guenon / Evola do not always find themselves in agreement with Tomberg. Nobody is trying to deny the intellect is vital, but one must also bear in the same mind that the Luciferian impulse acts directly on the mind. This is the problem with theosophy -- both Guenon and Evola appear to me as being basically theosophists from an intellectual point of view, albeit with catholic sentimentality, with all the attend avenues of illumination but lacking the purity of true Christian heart-centred faith. The problem with theosophy is well known but not necessarily understood -- hence the immense errors that early occultists fell into, in many cases unawares.

dear unknown friends, Christian Hermeticism has no pretentsions to rival either religion or official science. He who is searching here for `true religion' or `true philosopy', or `true science' is looking in the wrong direction. Christian Hermeticists are not masters, but servants. They do not have pretension (that is, in any case, somewhat puerile) of elevating themselvess above the holy faith of the faithful, or above the fruits of the admirable efforts of workers in science, or above the creations of artistic genius. Hermeticists are not guarding the secret of future discoveries in the sciences. They do not know, for example, just as everyone at present is ignorant of it, the effective remedy against cancer. Moreover they would be monsters if they were to guard the secret of this remedy against this bane of humanity without communicatin git. No, they do not know it, and they will be the first to recognise the superiority of the future benefactor of the human race, that savant who will discover this remedy.

likewise they recognise without reserve the superiority of a Francis of Assisi -- and of many others -- who was a man of the so-called `exoteric' faith. They know also that each sincere believer is potentially a Francis of Assisi. Men and women of faith, of science and of art are their superiors in many essential points. Hermeticists know it well and do not flatter themselves to be better, to believe better, to konw better or to be more competent... That which they possess does not comprise any tangible advantage or objective superiority with regard to religion, science and art; what they possess is only the communal soul of religion, science and art.

Letter 1, The Magician

Of course I guess this only applies if one sees one's self as a Christian Hermeticist first and foremost rather than a philosopher first and foremost...

\hfill

\texttt{Charlotte on 2012-05-25 at 10:27 said:}

Or Evola at any rate seems more theosophical... of the two authors I much prefer Guenon, he has the greater insight

\hfill

\texttt{logres on 2012-05-25 at 13:19 said:}

He's saying that there is no conflict -- I don't understand... ''no pretensions to rival... ''. I take that at face value. The primacy of the heart in union, but the primacy of the intellect in formally judging is the classic Catholic formulation -- St. John of the Cross himself helped formulate it. I don't see what your quarrel with that is.

\hfill

\texttt{Charlotte on 2012-05-26 at 08:43 said:}

err, I think you are forgetting that I was responding to a posting where `artists, writers' and so on were being given secondary place to `elect philosophers'... .there were also disagreements about who/what constituted the `elect' and I object to the majority of people being called atheist and perceived as lesser than unspecified `others'. I don't agree that those at the top of the pyramid are blameless while the `masses' are to blame for all the ills of the world. Where is the compassion in that?

Look, I've spent a lot of time on this, if nobody can understand what I'm trying to say then clearly it's energy not well spent! No hard feelings but it's probably best if I leave you to talk amongst yourselves :-)

\url{http://alchemical-weddings.com/alchemical-weddings/liberating-action}


\hfill

\texttt{logres on 2012-05-26 at 09:04 said:}

I understand perfectly what you are saying, \& agree with the information/wisdom contained in it. I just don't understand the contextual point of making it, except as a warning. Besides the ``masses'' you are speaking of are the ``folk'' and individuals, not the modern masses referred to by Guenon and others. Even bees have a hierarchy. It sounds like it's just a dirty word to you. That's fine, but it's not being used in that sense. Maybe we could just agree on another word.

\hfill

\texttt{logres on 2012-05-26 at 09:08 said:}

I brought up St. John of the Cross because he recognized the ``hierarchy'' of intellect in apprehending (practically) virtue. Thus, he didn't embrace a rhetoric of practical virtue or liberty, but instead achieved it (not that you are). Ditto with Tomberg. Let's agree to another word to denote ``hierarchy'' as understood in a non-spiritual sense, that is, a hierarchy which contradicts Scriptural exhortation to ``serve one another''. Perhaps, dominarchy or something like that? Are you familiar with Mouravieff's diagram?

\url{http://unurthed.com/2009/10/21/mouravieffs-correction/}

Hierarchy in THAT sense.

\hfill

\texttt{Graham on 2012-05-26 at 11:28 said:}

Charlotte said: ``it is to do with the primarcy of the heart over mind as our guide.''

Guenon and Evola agree with that forumulation. They say that man's symbolic centre is the heart, but they understand the heart as the seat of intellectual intuition united with will, while the mind is the seat of thought. And only in their sense can it be truly said that the heart guides the mind. You consistently give the impression of understanding the heart's Love (Caritas) as principally emotional; but in this case, love is blind. Emotion guiding thought is the blind leading the blind.

Another tradition Guenon points to is that Love and Knowledge are united, as warmth and luminosity are united in physical light. Luminosity guides (thus ``illuminate'' as ``to make known'') while warmth empowers.

This is why the fruit of Faith is knowledge of God, and why Faith `precedes' Love as a theological virtue. Even here, where Love is defined as willing the Good rather than in emotional terms, knowledge is the guide.

But Charlotte you do raise a good point. Cologero has demonstrated that Evola and Tomberg share a great deal -- perhaps a fundamental worldview -- as hermeticists. But how do we reconcile Evola's ``coldness'' and ``distance'' with Tomberg's ``warmth'' and ``nearness''? What is the source and significance of this difference?

\hfill

\texttt{Graham on 2012-05-26 at 11:55 said:}

Charlotte,

The artist engaged in his concrete function (painting, sculpting, and so on) is servile or imitative. Keep in mind that ``servile'' here is descriptive rather than pejorative. The artist, if and when he is held rapt in the contemplation of beauty, is at that moment a philosopher, not concretely an artist.

Now it's possible that a philosopher could be a quietist or an escapist, which I think is the phenomenon you object to. But if an artist can be a philosopher, then obviously a philosopher can be an artist, and in the most universal sense.

\hfill

\texttt{Will on 2012-05-26 at 20:31 said:}

The qualities listed as Transcendental in the chart are two of the three ``foods of the soul'' in Catholic philosophy -- the third being Beauty. The Good, the True and the Beautiful is a common triplet.

Beauty is related to love, which is the nature of God, according to the Gospels. Beauty / love / God is thus the unifying principle, as the article states.

This relates to the discussion of art, which conveys truth and goodness through beauty.

If we add this third to the chart, we might say that its principle is Unity, its faculty is Love, and its transcendental quality is Beauty.

We might also consider the correspondence of these three to the three Dan Tiens in Taoism, and also the three faculties of the soul in Platonism, though there isn't an exact correspondence. In Chinese thought, willing has its origin in the kidneys, which are the seat of the life-force. Love comes from the heart, and thought from the head. These are the bodily locations of the three Dan Tiens (``seas of elixir.'')

Platonism, in contrast, associates the lower torso (genitals and stomach) with appetite, the chest with will and courage, and the head with thought, contemplation and intuition. Love is still central to Platonism however, since the love of wisdom is the essence of the path, and for the philosopher is the ruler and harmonizer of the appetite, will and mind.

\hfill

\texttt{logres on 2012-05-26 at 22:52 said:}

Typo: ``not that you are'' I meant, not, you haven't apprehended virtue, but you haven't ``embraced rhetoric'' (over substance).

\hfill

\texttt{Charlotte on 2012-05-27 at 14:11 said:}

well thank you for your comments, I do believe you are sincere but I'm not sure how much more I can explain here of what I know to be true. I don't wish to argue or get overheated but when one is in a position of defender of faith it's difficult to just sit nodding quietly when one sees things that are contrary occuring in an environment that is professed to be -- and has the potential to be, with the right elements -- an active Hermetic community.

As far as I'm concerned this needs to be pure if I'm to play a part in it, but that isn't meaning to be judgmental to anyone here, I understand that there are other peopale with different types of interests and different reasons for looking into things. However my concerns regarding what is really needed in the spiritual dimension are sounding alarm bells very loudly at some of things I hear, at some of the spirit I sense, and I know I'm not the only person to have experienced such misgivings over Evola. Perhaps it would be helpful to ponder why I -- we -- have had reason to feel this way, even if it is not readily apparent to you?

So, I see some major blind spots with respect to doctrine and even understanding of the perennial philosophy (I certainly don't think Evola is the undisputed expert) but I'm also conscious that others here perceive the blind spot as being mine :-) But that is fine, we are all trying to deepen our knowledge and one of the ways in which we must achieve this is via mirrors... .

I see a dire conflict between some of the authors/philosophers being `upheld' and all together used -- as if they were of one mind, `mind' being all important, it seems -- to bolster both a particular elitist notion of tradition and also catholicism, which is bizarre given Evola's own religious views regarding Christianity. This is why I'm getting my knickers in a twist over terms being used, because I'm sensing some of the spirit that is underpinning them and I don't like it. I'm extremely careful regarding which `charged' works I read and I would be wary of Evola in any context. Guenon is a different story as I've mentioned a few times, but coldness and cruelty have no place WHATSOEVER in the great work of Christian Hermeticism. I really don't think Tomberg should be used to back up the ideas of Evola, it strikes me as wholly inappropriate, not least of all because Tomberg doesn't want rousing for debates like this, he doesn't like the attention at all. Politically Tomberg and Evola may have agreed, and doubtless they understood some of the same things, but I see no reason whatsover to persist in reviving Evola, those days are gone. Really those days are gone.

Which brings me onto another point relating to catholicism, which is more distressing and in fact it is probably best if I say nothing as it would give me no pleasure and you neither. I only allude to it so as to help you understand some of what has been troubling me about all this here... .bear in mind the bones of St Peter have already been sent (by the Pope) to Fatima. Don't be blinded by your own sentimentality with respect to personal religion to what is happening in this age. It gives me no joy to say this, in fact it gave me many tears for a number of years... Perhaps I'm repeating myself a bit now, but I also hear it said that intellectuality is something to be sniffed at (fair enough, I basically agree), and yet this all appears to me to be an intellectual endeavour and I'm not entirely sure what the overall agenda is, even though some interesting things are discussed and some interesting works brought to light. This in itself is not an unworthy task -- in fact it may be very worthwhile indeed, certainly some of the texts I've read as a consequence have been personally quite helpful -- but I sense there is confusion regarding `faith' that needs straightening out. The problem for the esotericist -- in particular a left hand path practitioner -- with respect to being led towards Islam, seemingly inevitably, are almost impossible to figure out unless you get back to the right track... .I know this from experience and have resolved it in a way I can fathom, but can all of you?

I realise Evola is not Christian and that Guenon had decided the only way was to go Sufi (I think, tho may have got this wrong?), which is their choice of course, but how does this fit in with catholicism? To the catholic sensibility it all seems very theosophic!! Now I am not a member of the exoteric catholic church so to me personally this isn't a great problem, but overall I see a lot of such contradictory threads that my logical mind is feeling exasperrated by that -- I think it was Rudolf Steiner who said that the logical person will feel actual pain when faced with illogicality, not that I'm wishing to nail the good doctor further into his coffin by mentioning him again :-)

Regarding 'emotion', yes I have a normal human sense of emotion... .it's also called passion, something quite close to the heart of the Christian faith, which is rightly pointed out to be about the spirit of service to mankind. Spirituality void of passion is a boat without sails, which doesn't mean to say we should ditch the rudder of reason. Evola seems to be lacking both passion and compassion, he veers to far into the dark side for my liking so if he's going to be the main focus I just can't do it anymore!!

\hfill

\texttt{Tony Ciapo on 2012-05-27 at 18:08 said:}

\textit{Dispassion the Aim of All Asceticism}\footnote{\url{http://orthodoxwayoflife.blogspot.com/2010/06/dispassion-aim-of-all-asceticism.html}}.

Unfortunately, Charlotte, the illogicality is arising from you. Clearly, we are picking and choosing from Evola, and there are reasons for that. First of all, many have become interested in Tradition because of him. We agree that in many ways, he has also distorted Tradition; that is something we have addressed many times. Yet, because he is not attached to any living tradition, we can use him to demonstrate various points without partisanship. It is part of a larger project and we will be moving beyond Evola.

Now Evola wrote an important book on Hermetism, which has been widely praised by the likes of Eliade and Jung. If you need to pass judgment, at least read that first.

As for the illogicality that distresses you, here are some examples:

1) You disagree with the idea of an elite, yet when you went to Turkey, you boasted about your meeting with a Sufi leader. Why didn't you discuss spirituality with the falafel vendor on the street corner?

2) You object to a certain characterization of the masses, one that Tomberg himself claimed to be true. As a matter of fact, it is historically true that the French and Russian revolutions --- ``in the name of the people'' --- were explicitly atheistic. That is all anyone is saying ... there is no disputing facts.

3) You claim the masses are ``naturally'' devout, but have been ``misled'' by some higher ups. That is precisely what we mean ... the masses are formless and are \emph{always and necessarily} formed by their superiors. So who should be the proper superiors? Obvisouly, a spiritual elite, in which case the devoutness and faith of the mass will be able to take form and thrive. I can't understand why any thoughtful person would object to this. Was not Jesus and the apostles an example of an elite group that spiritually transformed the West?

Please, Charlotte, before you nitpick and attack, please ask for clarification for things you don't understand. And please don't go off onto fantasies about the ``dark side'' and other such bizarre notions. Much of what we ``know to be true'' most certainly isn't. That is why every spiritual path begins with a process of purgation and catharsis, whereby our minds are cleared of our everyday beliefs.

Although we enjoy your lively comments, if what you find here makes you too uncomfortable, we won't take offense if you choose not to participate. But if you do decide to continue, please read again what Tomberg writes about what we have ``in common'', in our ``depths'', and not focus so much on personal experience, which is what divides, since it is contingent and changeable. This is found in Arcana IV.

\hfill

\texttt{logres on 2012-05-27 at 21:56 said:}

Charlotte, I think I may be as Christian to the bone as yourself, and certainly Evola was characterized by some as cold \& cruel -- I think one of his female acquaintances did in fact. Christ (however) had a special affinity for Roman soldiers (this was noted by Chesterton) and His sword has to be above justice or injustice -- ``those who are not against us are with us'' \& ``those who are not with us are against us'' -- He said both of these statements. Cologero has gone to pains to point out that Evola did not (actually) follow a living spiritual tradition, and that this fact makes him (if you wish to hear it plainly) ``less elite''. However, I can't help thinking (after each reading of his work) -- what a fantastic mind -- ``our Marcuse, only better''. Surely it would be part of the ``wisdom of the serpent'' to learn from him? There are dangers here, but Life is dangerous, regardless. Evola was a weapon resurrected from ashes, and it would be a service to the Faith to employ him at any and all points that his native worth demands the justice of doing so. And I do hope you continue to read and discuss here... Who else or where else can you gain a clinical diagnosis of the modern situation in such a short, accessible form? Your feelings do you credit, but they could be used against you -- I appeal to you to consider it and reconsider it before deciding.

\hfill

\texttt{logres on 2012-05-27 at 22:00 said:}

Mouravieff, Preface to Gnosis:

``Yet, paradoxically, many Europeans who feel drawn to these researches turn their eyes towards the non-Christian Traditions: Hinduism, Buddhism, Sufism and others. It would certainly be exciting to compare esoteric thinking in these different systems, because the Tradition is One, and whoever delves deeply into these studies will not fail to be struck by this essential unity. Yet to those who desire to go beyond pure speculation, the problem appears in a different light. This unique Tradition has been and still is now being presented in multiple forms, each meticulously adapted to the mentality and spirit of the human group to which its Word is addressed, and to the mission with which this group has been charged. For the Christian world, the easiest way, the least difficult way to reach the goal, is to follow the esoteric Doctrine which forms the basis of the Christian Tradition. Actually, the thought of a man who has been born and formed at the heart of our civilization, be he Christian or not, believer or atheist, is impregnated with twenty centuries of Christian culture. It is incomparably easier for him to begin his studies starting from this environment, rather than to adapt to the spirit of an environment different from his own. Transplantation is not without danger, and generally gives hybrid products.''

This last statement is the danger alluded to by Charlotte, but it's been addressed sufficiently by those operating the website to warrant the belief that they are proceeding with caution.

\hfill

\texttt{Tony Ciapo on 2012-05-27 at 23:52 said:}

Logres, lest you mislead casual readers, neither Evola, Guenon, nor anyone here has claimed to be part of any elite. What Evola did is describe his vision of how a \emph{future} elite may arise and manifest itself. As you yourself pointed out, this is not, strictly speaking, a merely human enterprise; to think so is to ignore the third dimension of history.

It is true that we prefer the Medieval civilization as the model for Tradition, for the reasons you specified. However, let me make it clear, there are pagan, Muslim, and Hindu readers who do not accept all the beliefs of that era. If you want to focus exclusively on ``Christian Hermetism'', I suggest you move the conversation to \url{http://www.MeditationsOnTheTarot.com}. I suggest the topic of ``freedom'', which Charlotte claims, following the Jacobins, is man's natural state. Jesus, OTOH, said only the Truth will make you free. The saints have claimed that man in his ``natural state'' is not free, but rather a slave to his passions. Which side do you plan to take, Logres? You can play Higgins to Pygmalion if that is your cup of tea, just don't fill up the comboxes here with your efforts. A further point is that we are not defending the Traditional worldview, but rather we are simply describing it as best we can. I'm sure it is quite shocking to someone wandering in the dark fog of the illusions of modernity. After all, isn't to be called ``medieval'' considered an insult in today's world?

Hence, while debate on specific points can be useful, we are not out to convert anyone to a specific worldview. In any case, only an intellectual conversion is of value, and that won't happen from reading a blog.

Moreover, critics are obliged to try to understand what is being said here and not resort to special pleading. For a guest to come here and accuse everyone else of ``illogicality'', and following the ``dark side'', on the basis of claims to special knowledge and even moral superiority is simply rude if not absurd.

\hfill

\texttt{Charlotte on 2012-05-28 at 06:46 said:}

Tony, I'm sorry if you thought I was being boastful -- I am not a boastful, prideful person, I'm actually very shy away from writing and have such terrible self esteem that it infuriates most people who know me personally. It really makes me wonder why I continue to get into these things, I must be a true masochist lol!

I only retell personal experiences when I think they might resonate with someone, or to show that I understand what they are trying to say. Unfortunatley many people interpret that as boastful, so most of the time I don't share them anymore. Certainly they don't appear to be of interest so are therefore not worthwhile. I don't believe I said the man in question was a Sufi `leader', although he might have been and I didn't know it. He was in fact a poor artist I met in a store while looking at pictures and admiring the work, and I actually thought people here might find the tone of the conversation amusing under the circumstances... it was all the more amusing from my point of view because at almost exactly the same time as I spoke to this Sufi I was harangued rather abusively on email by an extremely vehehment member of the Falun Gong who said he thought I wouldn't recognise Jesus if I bumped into him on the street, which of course had just happened, in name at least. I guess it's one of those situations where you just had to be there to get it, but you are right, this is probably not a good place to share personal experience because the concern here is with theory, and the medium of communication we're using is not ideal.

Too much gets lost in translation when people can't hear tone or read body language. Also, I didn't discuss spirituality with the falafel seller as I really didn't want to talk about spirituality with anyone and barely said a word to the Sufi either -- I didn't have anything I felt worthy to share, I only wanted to see Hagia Sophia.

And I wasn't accusing you all of following the dark side, I was saying I think there is a spirt of darkness coming from the Evola work that I'm sensitive to and I felt it my duty to point this out. Is that really an unfair assessment, bearing in mind I am new to his work? Do you only want to preach to the congregation or are you interested at all in introducing new people to the work of someone you believe is so valuable? Can you blame them for questioning what so many others have felt it necessary to question? I don't doubt for one minute that Cologero and others mean well, but I wouldn't be a very good friend if I did not sound a warning bell if and when I saw possible trouble. Why would you bother to resurrect Evola from ashes if you did not believe there was a spiritual force to his work that can act upon people. In me (as a new reader) you have a perfect guinea pig by which to gauge these reactions and yet you get angry with the creature when an effect occurs... .again, I'm missing the compassion that gives me cause to doubt the overall value, although I believe I have already pointed out that I think some of what Evola wrote is useful and that he is maybe even more helpful than Tomberg at one or two critical junturs pertaining to the labyrinth of the Kali Yuga. however, I'm not persuaded away from my basic view that the left hand path will only get you so far and if you want to escape the cycle of reincarnation (however this is to be understood), the only way out is the pure way.

I also felt that the purpose of the site was not clear, it seemed that on some levels an attempt was being made to found in the spiritual planes this `ideal hierarchy', for want of a better term -- a not unreasonable desire, such a thing is needed -- yet you now confirm that you're simply trying to describe the meaning of Evola et al. In the latter case then fair enough, there is no need for me to be worried. I am only concerned with the essence of the former (and maybe this concern is what's been affecting my assessment of the activities here, I can see that now) and the way it turns out to be, so in that case yes, I am better off focusing on the Meditations and prayer. I misunderstood your motives and intentions with this website and this is where my concerns have arisen from.

By the way I don't feel ashamed of my defence of humanity -- it's a tough job but someone has to do it. We all hate communism, but is the peasant really worse than a Stalin? And I'm not follow the `Jacobites', I'm not following anyone except Christ -- well, trying to anyway, even if I stumble and fall very often.

But maybe you should make this a closed forum if you don't want guests' opinions? I don't have to think you're converting people to have an opinion do I? Are any opinions welcomed at all if they differ at all from the overall consensus in any shape or form? One cannot get very far in the hermetic task of nailing down the opposites if one refuses to countenance any form of challenge to their ideas... And as for feelings and the value of them, well, I appreciate the warning that they'll be used against me as I understand it's made in a friendly spirit. The way I see it there's something called `balance', and it seems as if the balance here often veers so far towards that of what is seen as rationality (although I don't agree that everything presented as such here really is ultimate rationality) that one's desire for balance and temperace is greatly assaulted! There needs to be a counterweight and like I said before, if I'm the only one willing to play the fool then it won't be the first time... But Logres, thank you for your more kindly sentiments, they are certainly appreciated and I know you are a Christian. I'm not meaning to have digs at people that isn't my style at all -- in my mind there are very important principles at stake and I am willing to stand up for those, but I realise there are rocks and hard places lying around here and don't want myself or anyone else to get pointlessly bruised.

The work of catharsis and purification is -- according to what I know anyway -- achieved by focusing on Christ and love -- the mind can only take us so far, it's love alone that can bring us to our knees.

As for the relative value of being hot or cold, we can at least all take heart from the fact that we're not lukewarm!! blessings to you all Cx

\hfill

\texttt{Angolmois on 2012-05-28 at 11:27 said:}

It seems to me that Evola is attractive to a lot of masculine men who are of some ancient and even archaic kshatriya tendency, and this is very truthfully dangerous in todays spiritual climate. What Evola needs to be complemented with is Christian love and compassion, in theoretical and intellectual sense he is very much to the core in esoteric matters. Evola is too solar, I have said this many times, and it even caused his ``accident'' which happened in the course of the war and into the solar plexus chakra. We all need to learn from this: there is absolutely no room for spiritual violence after Christ.

There is a good text about intellect and love in Sophia Perennis website that discusses Guénon's and Schuon's possible prejudices against love:

\url{http://www.sophiaperennis.com/discussion-forums/general-discussion/were-rene-guenon-and-frithjof-schuon-biased-against-love/}


Charlotte: Guénon was a catholic when he lived ini France. He turned towards sufism because he moved into a muslim country. This is how traditional men and women act, it is part of being courteous and a way of ``choosing a tradition'' that fits ones intellectual and even aesthetic sensibilities.

I have seen a lot of criticism of Evola here, and that is a very good thing. You may be over-reacting, Charlotte, but I'm very glad that you have brought things up!

My blessings and love are with all of you, as always.

\hfill

\texttt{Matt on 2012-05-28 at 11:33 said:}

``By the way I don't feel ashamed of my defence of humanity''

Charlotte, I don't think anyone here is trying to make you feel ashamed for defending humanity. It seems to me that you are afraid that the idea of a spiritual elite (elect might be a better term) is synonymous with misanthropy,and its about getting revenge on the masses, which is the exact opposite of its intended meaning here. As previously stated, the elect is about renewal, one that is to the benefit of the mass (even if they can't see that benefit due to being caught up with the various illusions of modernity), which is reflective and therefore needs to be formed by that which is beyond it.

There is nothing wrong with being cautious... .whether it be about certain missteps in this or that thinker's thought, but too much caution can lead to a fear that is unhealthy.

\hfill

\texttt{Charlotte on 2012-05-28 at 13:58 said:}

Ah, thank you Angolmois, I was hoping I would manage to get my point across somehow and I'm glad it made sense on some level! yes, there is a lack of love in them, which is the single biggest problem currently afflicting humanity. With more love the world could be a far more beautiful, safe and magical place for everyone.

Also thank you for clarifying about Guenon, that makes some sense, although I have to wonder if it was simply an outer conversion? It is one thing to outwardly `be as the Romans when in Rome', but another to fully change one's inner religion?

Thank you for the article, I shall read this with interest. Giving preference to `Truth' over `Love' is a tenet of Eastern Buddhist based philosophy, with the Falun Gong I mentioned being a good example of this (from what I hear). This is one very popular worldview, of course, but it is not that of the Western mystery tradition, which makes a clear distinction between creator and creator, lover and Beloved. This is at the heart of the immense dangers of the `theosophic' philosophies of the kind promoted by the `new age movement', although it is very very difficult indeed to see this danger when one is in the thick of the magical output -- the blinding light of Lucifer is to be treated with more than caution -- better the cloud of unknowing in this part of the circle, or `mantle of the Virgin', as someone mentioned lately, the Armour of God. (LSD, by the way, long considered a product of the `hippy movement', was created by the Tavistock Institute for official purposes -- remember this is when satanism started to be spread in the western world ona mass scale. We should remember also how the illuminati principles (I won't say `elite') have influenced this movement -- this is actually rather scary if one thinks about it carefully).

Sure as anything we all need to put on some form of protection when traversing the zone of contradiction, belt of lies and realm of Maya and we all better believe that if we want to avoid getting deep fried! Maybe I'm overreacting... again it wouldn't be the first time... .buy maybe I'm not! Suffice it to say I felt sufficient cause to urge immediate assumption of the Helmet of Salvation for everyone, better safe than sorry -- we should all be fully aware of the types of works that have been used to stimulate the thought memes of the new age movement -- Blavatsky, Bailey and Crowley for sure -- what about Guenon? maybe he is not well known enough. Evola would have given them something to react against, in that respect the grey area can get so dark it's bound to work against the overall cause. you know even Teilhard to Chardin has been used to promote new age philosophy (Via Terrence McKenna). Of course the supreme danger of `information wars' is that there is a percentage of dazzling truth in there that lulls people into a sense of complacency with respect to the rest.

Matt, an accident with the solar plexus does not sound great at all (we've been speaking of occult wars remember, it's electrical information so sparks can be expected -- but full on melt downs are a scary business and it's better to keep watering those fuel rods!) -- indeed it is something to greatly fear and endeavour to avoid, the same go somehow damaging any of the chakras. Can you tell us more about how this incident occurred Angolmois?

\hfill

\texttt{Charlotte on 2012-05-28 at 14:37 said:}

I just read the start of an essay from someone in the UR group pointing out the immense dangers of inadvertently undergoing a counter initiation. I was actually going to write something like that earlier but as it was relating to personal experience and that's been banned I thought it best not to. Here's what `Arvo' has to say about it:

Those who attempt to overcome human limitations and yearn to attain knowledge and power need to be aware of the existence of what, using Rene Guenon's term, one may call the counter-initiation. They also need to have an idea of the various forms assumed by it, as well as of the different means it employs to achieve its goals. As a starting point, we may take the general idea that forces exist which attempt to insert themselves into the human domain, both at an individual and at a collective level, not only in order to mislead every aspiration to true spirituality, but also in order to create new currents, suggestions, and ideological systems. These are such as to blur the vision of truth, to falsify values, to favor the emergence of inferior influences, namely of every form of materialism, disorder, and subversion in civilizations. The opposition between the forces of ``good'' and ``evil'' is a commonplace in religions; however, this does not adequately deal with the issues at stake. It is not simply a question of the moral and religious order, but of an action much larger, objective, and concrete, which the representatives of religion themselves often fail to recognize, at times even falling unconsciously as victims to it.

See this is what I'm talking about people, I'm more convinced than ever I was right to flag warning signs in grey areas... 

\hfill

\texttt{logres on 2012-05-28 at 15:02 said:}

C, nevertheless, Tony \& Cologero are right to warn you about implicitly accepting the French Revolution. Man is born in chains, to himself. The Counter Initiation just makes use of that fact.

\hfill

\texttt{Charlotte on 2012-05-28 at 15:27 said:}

And how about this one:

Evola himself wrote that the aim of the ``chain'' of the UR Group, aside from ``awakening a higher force that might serve to help the singular work of every individual,'' was also to act ``on the type of psychic body that begged for creation, and by evocation to connect it with a genuine influence from above,'' so that ``one may perhaps have the possi\-bility of working behind the scenes in order to ultimately exert an effect on the prevailing forces in the general environment.''

See how can it be wrong to double-check in this situation?? I knew it was there, impossible to hide, I was just taking part in a chain reaction :-)

\end{sffamily}
\end{footnotesize}

